\section{Resultados e Discussão}
Esta seção apresenta os resultados obtidos e a análise crítica dos achados. Podem ser incluídos gráficos, tabelas, quadros, equações e comparações com trabalhos relacionados. Caso a instituição prefira, esta seção pode ser dividida em duas: Resultados e Discussão.

\subsection{Apresentação dos Resultados}
Descreva os principais resultados obtidos no desenvolvimento do trabalho, destacando dados, evidências, medições, protótipos, experimentos ou demais entregáveis relacionados aos objetivos da pesquisa.

\subsection{Discussão dos Achados}
Analise criticamente os resultados, comparando-os com a fundamentação teórica e com trabalhos relacionados. Indique interpretações, limitações, implicações práticas e possíveis causas para os resultados encontrados.

Esta seção apresenta orientações práticas inspiradas no modelo \texttt{conference\_101719.tex}, adaptadas para o contexto de um Trabalho de Conclusão de Curso. O objetivo é explicar os principais recursos do \LaTeX{} usados neste template e indicar como manter a padronização visual do documento.

\subsection{Estrutura do Documento}
O arquivo principal, \texttt{main.tex}, organiza o trabalho por meio de comandos \verb|\input|. Essa estratégia permite separar o conteúdo em elementos pré-textuais, textuais e pós-textuais, facilitando a manutenção do projeto. Assim, cada seção pode ser editada em um arquivo próprio sem alterar a estrutura geral do documento.

A classe \texttt{unisatex} é responsável por carregar a base tipográfica e estrutural derivada do \texttt{IEEEtran}. Ela também disponibiliza comandos próprios do modelo institucional, como \verb|\unisatexsubtitle|, \verb|\unisatexauthorblockN| e \verb|\unisatexauthorblockA|.

\subsection{Título, Subtítulo e Autores}
O título do trabalho é definido com \verb|\title|. Quando houver subtítulo, recomenda-se usar \verb|\unisatexsubtitle| para manter a apresentação padronizada:

\begin{verbatim}
\title{Título do TCC\\
  \unisatexsubtitle{Subtítulo}}
\end{verbatim}

Os autores devem ser definidos dentro de \verb|\author|. Para manter a compatibilidade com o formato de conferência, cada autor utiliza um bloco de nome e um bloco de afiliação. O comando \verb|\thanks| pode ser empregado para notas de rodapé, como a indicação do curso do aluno ou a função do orientador.

\subsection{Resumo, Abstract e Palavras-chave}
O resumo em português utiliza o ambiente \verb|abstract|, enquanto as palavras-chave utilizam o ambiente \verb|keywords|. Para a versão em inglês, recomenda-se envolver o conteúdo em \verb|otherlanguage| e usar o ambiente \verb|englishkeywords| para que o rótulo seja apresentado como \textit{Keywords}.

Esses elementos devem ser objetivos e não numerados, pois pertencem aos elementos pré-textuais do TCC. O resumo deve apresentar problema, objetivo, metodologia, resultados e conclusão de forma sintética.

\subsection{Seções e Subseções}
As divisões textuais são criadas com \verb|\section| e \verb|\subsection|. O \LaTeX{} numera automaticamente as seções, mantendo a consistência do documento. Evite numeração manual nos títulos, pois isso prejudica referências cruzadas e futuras reorganizações do texto.

\subsection{Citações e Referências com BibTeX}
As citações devem ser feitas com \verb|\cite|, como em \verb|\cite{ieeetran}|. As referências bibliográficas ficam no arquivo \texttt{references.bib}, e a lista final é gerada automaticamente por BibTeX a partir dos comandos:

\begin{verbatim}
\bibliographystyle{IEEEtran}
\bibliography{references}
\end{verbatim}

Essa abordagem reduz erros de formatação e garante que apenas obras citadas no texto apareçam na lista de referências.

\subsection{Equações}
Equações devem ser inseridas com o ambiente \verb|equation| quando precisarem de numeração. Para citar uma equação no texto, recomenda-se usar \verb|\label| e \verb|\eqref|, evitando referências manuais que podem ficar incorretas caso a ordem do documento seja alterada.

\begin{equation}
a + b = \gamma\label{eq:exemplo-latex}
\end{equation}

A Equação~\eqref{eq:exemplo-latex} ilustra uma equação numerada automaticamente.

\subsection{Figuras e Tabelas}
Figuras e tabelas devem ser inseridas próximas ao trecho em que são citadas. Em geral, tabelas usam legenda acima, enquanto figuras usam legenda abaixo. Também é recomendável usar \verb|\label| depois de \verb|\caption| para permitir referências automáticas.

\begin{table}[htbp]
\caption{Exemplo de tabela no modelo UniSATeX}
\begin{center}
\begin{tabular}{|c|c|}
\hline
\textbf{Elemento} & \textbf{Função} \\
\hline
\texttt{main.tex} & Organiza os arquivos do TCC \\
\hline
\texttt{references.bib} & Armazena referências BibTeX \\
\hline
\end{tabular}
\label{tab:exemplo-latex}
\end{center}
\end{table}

A Tabela~\ref{tab:exemplo-latex} mostra um exemplo de organização de informações em duas colunas.

\begin{figure}[htbp]
\centerline{\includegraphics{assets/fig1.png}}
\caption{Exemplo de imagem inserida no modelo UniSATeX.}
\label{fig:exemplo-imagem}
\end{figure}

A Figura~\ref{fig:exemplo-imagem} demonstra como inserir uma imagem armazenada na pasta \texttt{assets}. O caminho deve ser relativo ao arquivo principal do projeto.

\subsection{Boas Práticas de Escrita e Formatação}
Antes de ajustar detalhes visuais, recomenda-se escrever e revisar o conteúdo textual. O modelo já define margens, colunas, fontes, espaçamentos e estilo de referências; portanto, alterações manuais nesses itens devem ser evitadas. Também é importante remover textos explicativos do template antes da entrega final do TCC.